\documentclass[kulak]{kulakarticle} % options: kulak (default) or kul

\usepackage[dutch]{babel}
\graphicspath{{afbeeldingen/}}

\newcommand{\vscode}{\includegraphics[width=1.5ex]{vscode}\texttt{VS}\hspace{.1ex}\texttt{Code}}
\newcommand{\git}{\includegraphics[width=1.5ex]{git}\texttt{Git}}
\newcommand{\github}{\includegraphics[width=1.5ex]{github}\texttt{GitHub}}

\title{Werken met \git\ en \github\ vanuit \vscode}
\author{Stijn Rebry}
\date{Academiejaar 2026 -- 2027}
\address{
	\textbf{Groep Wetenschap \& Technologie Kulak} \\
	Teamopdracht}

\begin{document}
	
	\maketitle
	
	\tableofcontents
	
	\section*{Inleiding}
	
	Het programma \git\ en de webinterface \github\ zijn hulpmiddelen om versies van bestanden bij te houden en met meerdere personen samen te werken aan een project. Oorspronkelijk werden ze ontwikkeld voor softwareontwikkeling, maar ondertussen worden ze ook gebruikt voor verslagen, onderzoeksprojecten, websites, presentaties, data-analyse en vele andere toepassingen waarbij meerdere versies van bestanden ontstaan.
	
	De basisidee is eenvoudig. In plaats van telkens bestanden op te slaan onder een andere naam zoals \texttt{versie1}, \texttt{versie2} of \texttt{eindversie\_definitief\_echt\_definitief}, bewaart \git\ een gedocumenteerde geschiedenis van alle wijzigingen. Op elk moment kan worden nagegaan wie een wijziging heeft uitgevoerd, wanneer die wijziging gebeurde en hoe een eerdere versie van een bestand eruitzag.
	
	\git\ bestaat essentieel uit honderden tekstcommando's die via de commandoregel kunnen worden uitgevoerd en biedt talloze geavanceerde mogelijkheden voor versiebeheer, samenwerking, testen en automatisering. In de praktijk werken veel gebruikers niet rechtstreeks met deze commando's, maar via een grafische gebruikersinterface die de meest gebruikte \git-functies toegankelijk maakt met knoppen, menu's en overzichtsvensters.
	
	\github\ is een online platform dat \git-repositories kan bewaren en delen. Het voegt bovendien extra functionaliteit toe voor samenwerking, zoals issues, taakverdeling en projectbeheer. \github\ zelf vervangt \git\ niet, maar bouwt er bovenop voort.
	
	\git\ kan vanuit verschillende programma's worden gebruikt. In deze handleiding kiezen we voor één werkwijze: \git\ aansturen vanuit Visual Studio Code (\vscode). Waar mogelijk gebruiken we de grafische interface van \vscode\ en \github. Slechts af en toe zullen we een opdracht uit het \texttt{Command Palette} gebruiken. Daardoor kunnen we ons concentreren op de onderliggende ideeën zonder eerst een lange lijst \git-commando's te moeten leren.
	
	Na het doorlopen van de installatie-intructies, starten we met introduceren van versiebeheer op de eigen lokale computer (local repository, commit -- sectie \ref{sec:local}), om die gegevens vervolgens te publiceren op \github\ (remote repository, push, pull -- sectie \ref{sec:remote}). Daarna onderzoeken we hoe verschillende versies van hetzelfde project naast elkaar kunnen bestaan (branches, fetch -- sectie \ref{sec:branches}) en hoe deze versies daarna opnieuw kunnen worden gecombineerd (\emph{merge} -- sectie \ref{sec:merge}). In een groter project waar meerderen samenwerken, is het belangrijk het overzicht te bewaren van wat er nog te gebeuren staat en wie wat doet (\emph{issues}, \emph{assignees} -- sectie \ref{sec:issues}). Ten slotte passen we deze workflow toe op de teamopdracht van het vak (sectie \ref{sec:teamopdracht}).
	
	Na het doorwerken van deze handleiding beheers je de basisworkflow die ook in professionele software- en onderzoeksprojecten wordt gebruikt, om de bijdragen van verschillende personen gecontroleerd samen te voegen.

	\begin{figure}
		\includegraphics[width=.99\textwidth]{VSCode-schermonderdelen.pdf}
		\caption{Benaming van de belangrijkste schermonderdelen in \vscode\ met betrekking tot \git}
	\end{figure}

	\section{Software installeren, account activeren}
	\label{sec:init}

	\subsection{\vscode\ installeren}

	\begin{itemize}
	\item Download de laatste versie van \vscode\ via \href{https://code.visualstudio.com}{\texttt{code.visualstudio.com}}
	\item Selecteer de gewenste taal, aanvaard de gebruiksovereenkomst en klik door de volgende dialoogschermen: tenzij je andere voorkeuren hebt, zijn de standaardinstellingen goed.
	\item Klik op \texttt{Install}.
	\end{itemize}

	\subsection{\github\ account maken}
	\begin{itemize}
	\item Ga naar de webpagina \href{https://github.com/}{\texttt{github.com}} om een account te maken.
	\item Vul je KU Leuven-e-mailadres in en klik op \texttt{Sign up for GitHub}.
	\item Kies een veilig paswoord, hergebruik geen paswoorden.
	\item Kies als username hetzelfde als je e-mailadres, maar zonder symbolen \texttt{.} of \texttt{@} maar inclusief het eventuele cijfer \texttt{1} of \texttt{2} achter je naam. Laat het laatste gedeelte weg als dat de maximale lengte van een username overschrijdt (39 karakters)\footnote{Dit is misschien overkill... hoe?}. Dit is belangrijk opdat we jullie accounts automatisch zouden kunnen terugvinden.
	\item Klik op \texttt{Create account}.
	\item Je krijgt een e-mail met een beveiligingscode die je moet invullen op de volgende pagina. Daarna is je account actief.
	\end{itemize}

	\subsection{Account koppelen aan \github-Education}

	\begin{itemize}
	\item Log in op \github.
	\item Open de pagina \href{https://github.com/settings/education/}{\texttt{github.com/settings/education/benefits}} en klik op \texttt{Start an application}.\footnote{Is dit nodig?}
	\item Op basis van je e-mailadres herkent \github\ jouw instelling als Catholic University of Leuven, klik \texttt{Select this school}, \texttt{Share location} en \texttt{Continue}.
	\item Stuur een foto van je studentenkaart om toegang te vragen en volg eventuele verdere instructies.
	\end{itemize}

\subsection{\git\ installeren}
\begin{itemize}
\item Windows-gebruikers die niet eerder \git\ gebruikten, moeten dit programma eerst nog installeren. Download het installatieprogramma op \href{https://git-scm.com/}{\texttt{git-scm.com}}.
\item Start de installatie, aanvaard de gebruiksovereenkomst en klik door de volgende dialoogschermen: tenzij je andere voorkeuren hebt, zijn de standaardinstellingen goed.
\item Tot slot is het nodig om je gebruikersinformatie te configureren. Open de terminal in \vscode\ (\texttt{ctrl+`}) en voer daarin volgende commando's uit:
\begin{verbatim}
git config --global user.name "Voornaam Naam"
git config --global user.email "jouw-e-mail@student.kuleuven.be"
\end{verbatim}
\end{itemize}


	\section{Local repository, commits}
	\label{sec:local}
	
	\subsection{Lokale map maken} \git-projecten zullen we opslaan in een lokale map, \emph{niet} in de cloud (zoals OneDrive, Google Drive, Dropbox, ...). \git\ schrijft immers voortdurend bestanden weg en services als OneDrive zullen diezelfde bestanden tegelijk synchroniseren, wat conflicten, corrupte bestanden of prestatieproblemen kan veroorzaken.
	\begin{itemize}
		\item Open Windows Explorer.
		\item Typ in de adresbalk ``\texttt{C:\textbackslash}''.
		\item Klik links in de knoppenbalk op \texttt{Nieuw}, \texttt{Map}.
		\item Geef de map de naam ``\texttt{Projecten}''.
		\item Maak nog een map aan in \texttt{C:\textbackslash Projecten} met de naam ``\texttt{gittest}''.
	\end{itemize}
	Het hele punt van \git\ is natuurlijk wel om ook een kopie van de repository in de cloud te bewaren, als backup en om te delen met anderen, maar dat regelen we later in deze handleiding.
	
	\subsection{Lokale map openen in \vscode} Een project bestaat typisch niet uit één bestand, maar uit heel wat bestanden die worden verzameld in een map. We openen in \vscode\ de lokale map die net is gemaakt.
	\begin{itemize}
		\item Als je al andere bestanden hebt geopend in \vscode, open dan een nieuw venster, \texttt{File}, \texttt{New Window} (\texttt{ctrl+shift+N}).
		\item Open de map \texttt{C:\textbackslash Projecten\textbackslash gittest} in \vscode\ via \texttt{File}, \texttt{Open Folder}.
		\item Maak meteen een nieuw tekstbestand aan (\texttt{ctrl+N}): \texttt{File}, \texttt{New Text File}. Typ hierin enkel je naam (lijn 1) en opleiding (lijn 2).
		\item Sla dit bestand op (\texttt{ctrl+S}): \texttt{File}, \texttt{Save} met als naam ``\texttt{gegevens.txt}''.
	\end{itemize}
	
	\subsection{Lokaal repository initialiseren} 
	Op dit moment is de werkmap een gewone bestandsmap zonder versiecontrole. Om die te activeren, maken we er een \git-repository van.
	\begin{itemize}
		\item Klik in \vscode\ op het icoon \texttt{Source Control} (links in de knoppenbalk, \texttt{ctrl+shift+G}).
		\item Klik op \texttt{Manage Workspace Trust} en daarna op \texttt{Trust}, sluit het venster.
		\begin{itemize}
			\item \texttt{Trust}: \vscode\ beschouwt de map als veilig en schakelt alle functionaliteit in, waaronder \git-integratie, extensies en het uitvoeren van code.
			\item \texttt{Don't Trust}: \vscode\ opent de map in \texttt{Restricted Mode}. Bestanden kunnen nog worden bekeken en aangepast, maar bepaalde functionaliteit wordt uit veiligheidsoverwegingen uitgeschakeld.
		\end{itemize}
		\item Klik op \texttt{Initialize Repository} van de werkmap een repository te maken.
		\begin{itemize}
			\item Je ziet dat het bestand \texttt{gegevens.txt} wordt vermeld onder \texttt{Changes}. Onderaan in de zijbalk is de \texttt{Graph} nog leeg, hier verschijnt later de inhoud van de repository.
		\end{itemize}
	\end{itemize}
	
	\subsection{Eerste commit} Nu kunnen we een toevoegen van het eerste bestand aan de repository doorvoeren. 
	\begin{itemize}
		\item Tik in het vak \texttt{Message} (onder \texttt{Changes} en boven de \texttt{Commit}-knop) een beschrijving van de aanpassingen, in dit geval bijvoorbeeld ``Initiële inhoud''.
		\item Klik op \texttt{Commit}. Een pop-up vraagt ``stage all changes and commit?'', kies voor \texttt{Yes}. De rol van stage komt aan bod in sectie \ref{ssec:stage}.
		\item Klik op \texttt{Initiële inhoud} onder \texttt{Graph}.
		\begin{itemize}
			\item De \texttt{gegevens.txt} verschijnt met daarnaast de letter \texttt{A} (added file).
		\end{itemize}
	\end{itemize}
	
	\subsection{Aanpassingen doorvoeren aan een bestand} Nu zullen we het bestand \texttt{gegevens.txt} aanpassen en de aanpassingen doorvoeren naar de repository.
	\begin{itemize}
		\item Voeg aan het bestand \texttt{gegevens.txt} een derde lijn toe met daarop \texttt{Woonplaats: Gallifrey} (dat is wellicht fout, maar corrigeren we later). Sla het bestand op (\texttt{ctrl+S}).
		\item De bestandsnaam \texttt{gegevens.txt} is nu \emph{drie} keer te zien in de \vscode-omgeving:    
		\begin{itemize}
			\item Het bestand in het editorvenster, met de laatste aanpassingen, is opgeslagen maar zit nog niet in de repository.
			\item Klik op \texttt{gegevens.txt} onder \texttt{Changes}. Een nieuw venster opent in de editor met daarin de verschillen tussen de recentste versie van het bestand en de versie die in de repository staat. De toegevoegde lijn wordt groen gemarkeerd.
			\item Klik op \texttt{gegevens.txt} onder \texttt{Graph}. Een nieuw venster opent in de editor met daarin de versie van het bestand zoals die is opgeslagen in de meest recente commit van de repository. Deze bevat nog maar twee lijnen.
		\end{itemize}
		\git\ bewaart dus zowel de huidige bestanden in je werkmap als eerdere versies in de repository.
		\item Typ opnieuw een beschrijving van de aanpassingen in het vak \texttt{Message} onder \texttt{Changes}, bijvoorbeeld ``Woonplaats toegevoegd'', klik op \texttt{Commit}, klik op \texttt{Yes} in de pop-up. 
		\begin{itemize}
			\item De \texttt{Graph} toont nu twee commits, met de meest recente bovenaan en telkens met het bestand \texttt{gegevens.txt}. Door er op te klikken kan je de versies vergelijken.
			\item Naast de meest recente commit staat de letter \texttt{M} (modified file).
		\end{itemize}
		Met een commit leg je een versie van je project permanent vast in de repository, de message laat toe om later gericht een bepaalde versie terug te vinden. De repository bevat dus alle versies van je project, ook al zijn bestanden aangepast of zelfs verwijderd uit de werkmap.
	\end{itemize}
	
	\subsection{Bestanden toevoegen en verwijderen} Nu zullen we nieuwe bestanden toevoegen en een bestaand bestand verwijderen.
	\begin{itemize}
		\item Maak twee nieuwe bestanden (\texttt{ctrl+N}) met als namen \texttt{brol.txt} (willekeurige inhoud) en \texttt{steden.txt} (bevat vier lijnen met daarop respectievelijk \texttt{Brugge}, \texttt{Kortrijk}, \texttt{Leuven} en \texttt{Poelkapelle}). Sla deze bestanden op (\texttt{ctrl+S}) en voer een commit uit, klik op \texttt{Yes} in de pop-up.
		\begin{itemize}
			\item In de \texttt{Graph} staan inmiddels drie commits, met daarbij telkens enkel de bestanden die in die commit zijn toegevoegd (\texttt{A}) of aangepast (\texttt{M}).
		\end{itemize}
		\item Verwijder via \texttt{Explorer} (\texttt{ctrl+shift+E}) het bestand \texttt{brol.txt} uit de map (selecteren + \texttt{Delete}), bevestig met \texttt{Move to recycle bin}.
		\item Schrap het gehucht \texttt{Poelkapelle} uit de opsomming in \texttt{steden.txt}, sla het bestand op en commit de verwijderingen, klik op \texttt{Yes} in de pop-up.
		\begin{itemize}
			\item In \texttt{Source Control} wordt het bestand \texttt{brol.txt} gemarkeerd met de letter \texttt{D} (deleted file).
			\item De geschrapte lijn staat rood gemarkeerd in het bestand \texttt{bestanden.txt} onder de recentste commit in de \texttt{Graph}.
		\end{itemize}
	\end{itemize}
	Ook al is de informatie uit de map of uit een bestand verwijderd (zie \texttt{Explorer}), toch verdwijnt er geen informatie uit repository (zie \texttt{Source Control}).
	
	\subsection{Stage -- selecteren van aanpassingen voor commit}
	\label{ssec:stage}
	Hierboven werd een stap overgeslagen bij het committen. In een realistische workflow komt tussen opslaan en committen van aanpassingen nog het stagen. Niet alle wijzigingen die tegelijk in de werkmap aanwezig zijn, horen immers noodzakelijk in dezelfde commit thuis. Bekijk volgend voorbeeld: bij het schrijven van het eindverslag wordt figuur gemaakt en geïnterpreteerd in de tekst.
	\begin{itemize}
		\item Het maken van de figuur, verzorgen van de assen en de legende gebeurt in een \texttt{Python}-notebook waar alle geschreven code wordt getest en geïllustreerd.
		\item Het invoegen van de figuur, gebeurt in een \LaTeX-file. Er wordt een bijschrift gemaakt dat toelichting geeft bij de verschillende elementen op de figuur en een vewijzing in de tekst geschreven naar deze figuur.
		\item Tijdens het werken worden beide bestanden afwisselend aangepast: de figuur wordt verbeterd, daarna wordt het bijschrift aangepast, vervolgens wordt opnieuw aan de figuur gewerkt, enzovoort.
		\item Deze aanpassingen kunnen dan in twee commits worden doorgevoerd in de repository: één voor alles wat met het Python-notebook gebeurt en een tweede voor de aanvullingen in het verslag.
	\end{itemize}

	Hiervoor gebruikt \git\ de stage, wat leidt tot de volgende meer geavanceerde workflow Save -- Stage -- Commit. In het vervolg van deze tekst zullen we steeds alle aanpassingen tegelijk stagen, klik in de pop-up dus steeds op \texttt{Yes}, of eenmalig op \texttt{Always}.

		\subsection{Terminologie} Hiermee zijn alvast drie belangrijke begrippen geïntroduceerd:
		\begin{itemize}
			\item \emph{Werkmap}: de lokale map waarin de bestanden van het project staan.
			\item \emph{Stage}: selectie van wijzigingen die in de volgende commit zullen terechtkomen.
			\item \emph{Commit}: een opgeslagen versie van de werkmap op één bepaald moment.
			\item \emph{Repository}: het geheel van alle commits en dus alle aanpassingen en versies van elk bestanden.
		\end{itemize}

	\section{Repository publiceren op \github}
	\label{sec:remote}
	
	Tot nu toe staat de repository alleen op de eigen computer. Dat is niet erg veilig, je kan informatie verliezen door een defect of diefstal van je laptop. Het is op die manier ook niet mogelijk om samen te werken met anderen. Daarom zullen we een kopie van de repository online bewaren, op \github.
		
	\subsection{Repository publiceren vanuit \vscode} De lokale repository \texttt{gittest} zullen we nu publiceren op \github.
	\begin{itemize}
		\item Klik in \texttt{Source Control} op \texttt{Publish Branch}. Deze knop staat er enkel als er geen hangende aanpassingen meer zijn die op een commit wachten.
		\item Vervolg het proces om in te loggen met de account die gelinkt is aan je KU Leuven-e-mailadres.
		\item Kies \texttt{Publish to github Private Repository}..
		\begin{itemize}
		\item Een private repository is enkel zichtbaar voor de eigenaar en voor personen die expliciet toegang krijgen.
		\item Een public repository is zichtbaar voor iedereen.
		\end{itemize}
		\item Meld eventueel aan met je \github-account.
		\item Bevestig de publicatie.
	\end{itemize}
	
	\subsection{Repository bekijken op \github}
	Na enkele seconden verschijnt de repository ook in de browser op de \github-website.
	\begin{itemize}
		\item In de browser op je \github-pagina klik je door naar de repository \texttt{gittest}.
		\item Verifieer dat de bestanden \texttt{gegevens.txt} en \texttt{steden.txt} zichtbaar zijn, met inhoud volgens de recentste commit.
	\end{itemize}
	Op dit moment bevat \github\ dus dezelfde commits als de lokale repository.
	
	\subsection{Lokale wijzigingen maken}
	We zullen nu een nieuwe wijziging maken aan de lokale repository, in \vscode.
	\begin{itemize}
		\item Voeg aan \texttt{gegevens.txt} een vierde lijn toe met je favoriete kleur en sla het bestand op (\texttt{Ctrl+S}).
		\item Klik op \texttt{Commit}, verifieer de wijzigingen in \texttt{explorer} en vernieuw het bestand \texttt{gegevens.txt} op \github\ (\texttt{F5}).
		\begin{itemize}
			\item Zoals verwacht is het lokale bestand aangepast.
			\item De favoriete kleur is echter nog niet zichtbaar op \github.
			\item In \texttt{Source Control} staat de laatste commit in een andere kleur dan de voorgaande commits. Twee icoontjes (concentrische cirkels en een wolkje) geven aan welke de actieve lokale commit is, en welke de recentste commit is die op \github\ staat.
		\end{itemize}
	\end{itemize}
	De commit is lokaal opgeslagen, maar nog niet naar \github\ gestuurd.
	
	\subsection{Push uitvoeren} Om de nieuwe commit naar \github\ te sturen, gebruiken we een \emph{push}.
	\begin{itemize}
		\item Klik in \texttt{Source Control} op \texttt{Sync Changes}.
		\item Vernieuw de \github-webpagina (\texttt{F5}) en verifieer dat de kleur nu wel zichtbaar is in \texttt{gegevens.txt}.
	\end{itemize}
	
	\subsection{Wijzigingen rechtstreeks op \github\ aanbrengen} Niet alle wijzigingen moeten vanaf de eigen computer gebeuren. Een bestand kan ook rechtstreeks op \github\ worden aangepast (of, later, door iemand anders).
	\begin{itemize}
		\item Open op \github\ het bestand \texttt{gegevens.txt} en klik op het potlood-icoon (\texttt{Edit this file}). \item Zet op een nieuwe lijn je favoriete hobby.
		\item Klik op \texttt{Commit changes} in \github\ en bevestig.
		\item Ga na of de hobby in \vscode\ zichtbaar is. Dat blijkt nog niet het geval te zijn.
	\end{itemize}
	
	\subsection{Pull uitvoeren} 
	\github\ bevat nu een recentere versie dan de lokale repository, om de wijzigingen van \github\ naar de eigen computer te halen, gebruiken we een \emph{pull}.
	\begin{itemize}
		\item Ga naar \texttt{Source Control} (\texttt{ctrl+shift+G}).
		\item Klik op het icoontje \texttt{Pull} naast  \texttt{Graph}.
		\item Verifieer dat de hobby nu ook lokaal zichtbaar is.
	\end{itemize}
	
	\subsection{Terminologie}
	\begin{itemize}
		\item \emph{Local repository}: staat op je eigen computer.
		\item \emph{Commit}: wijzigingen vastleggen in de lokale repository.
		\item \emph{Remote repository}: staat online, op \github.
		\item \emph{Push}: lokale wijzigingen naar \github\ sturen.
		\item \emph{Pull}: wijzigingen van \github\ naar je eigen computer halen.
	\end{itemize}
	
	
	\section{Werken met branches}
	\label{sec:branches}
	
	Tot nu toe hebben we telkens rechtstreeks gewerkt op de hoofdversie van de repository. Soms willen we echter veranderingen uitproberen zonder de bestaande versie te wijzigen. Denk aan een computerprogramma waarvan versie 1 correct werkt, maar waarvan je een deel wil herschrijven om het sneller te maken. Zolang je de nieuwe versie nog niet getest hebt, en geverifieerd dat deze effectief sneller is, wil je de bestaande versie niet overschrijven maar behouden en blijven gebruiken.
	
	Daarvoor gebruiken we een \emph{branch}. Een branch is een alternatieve versie van dezelfde repository. Elke branch heeft zijn eigen geschiedenis van commits.
	
	\subsection{De huidige branch bekijken}
	
	\begin{itemize}
		\item Helemaal onderaan links in \vscode\ staat de naam van de huidige branch, op dit moment heet deze ``\texttt{main}''. Dit is de naam van de hoofdversie van de repository. Tot nu toe werden alle commits op deze branch gemaakt.
		\item Verifieer dat dezelfde branchnaam op \github\ vermeld staat. De branch wordt er met hetzelfde symbool aangeduid als de \texttt{Source control}-knop in \vscode.
	\end{itemize}
	
	\subsection{Een nieuwe branch maken} We zullen nu een nieuwe branch maken waarin we wijzigingen kunnen uitproberen.
	
	\begin{itemize}
		\item Klik in \vscode\ onderaan links op de branchnaam \texttt{main}.
		\begin{itemize}
			\item Er verschijnt een optie-lijstje onder de \texttt{Command Palette} met mogelijke commando's en een overzicht van de bestaande branches, momenteel enkel de branch \texttt{main} in de lokale repository en de (momenteel identieke) branch \texttt{origin/main} in de remote repository.
		\end{itemize}
		\item Kies voor \texttt{Create New Branch}, typ de naam ``\texttt{hobbies}'' in het tekstveld en druk op \texttt{Enter}.
		\item Door opnieuw links onderaan op de branchnaam te klikken, kan je snel wisselen tussen \texttt{hobbies} en \texttt{main}. Verifieer dat beide branches momenteel nog identiek zijn.
		\item Verifieer dat de nieuwe branch \texttt{hobbies} momenteel enkel lokaal bestaat en nog niet te zien op \github.
	\end{itemize}
	
	\subsection{Een wijziging maken op de nieuwe branch, local repository} We zullen nu een wijziging maken op de branch \texttt{hobbies}.
	
	\begin{itemize}
		\item Selecteer de branch \texttt{hobbies} en maak daarin een nieuw bestand \texttt{hobbies.txt}.
		\item Zet in het bestand een overzicht van je hobbies.
		\item Voer een \texttt{Commit} uit.
		\item Standaard worden in \texttt{Graph} enkel de commits van de huidige branch getoond. Klik naast \texttt{Graph} op het \texttt{Source Control}-icoontje, selecteer \texttt{All} en klik op OK.
		\item Schakel heen en weer tussen beide branches door op de branchnaam in de statusbalk te klikken.
		\begin{itemize}
			\item Het nieuwe bestand is enkel zichtbaar is in de branch \texttt{hobbies}.
			\item De \texttt{Graph} toont de geschiedenis van de verschillende branches met verschillende kleuren. Labels en icoontjes geven aan waar elke branch zich momenteel bevindt. Het wolkje geeft de meest recente bekende versie van de (\texttt{main}-)branch op \github\ aan.
			\item De inhoud van de werkmap in het venster \texttt{Explorer} en ook in het besturingssysteem zelf, verandert naargelang in \vscode\ een andere branch wordt geselecteerd.
			\item De geschiedenis van commits, branches en andere repositorygegevens wordt bewaard in de verborgen submap \texttt{.git}.
		\end{itemize}
	\end{itemize}
	
	\subsection{Een nieuwe branch op de \github-repository} Nu maken we een nieuwe branch op \github\ aan en voegen er ook een extra bestand aan toe.
	
	\begin{itemize}
		\item Open de repository in \github\, naast de branch-naam \texttt{main} staat de tekst ``1 Branch''. Klik hierop en daarna op \texttt{New Branch}. Geef de branch de naam ``\texttt{vakken}''. In het veld \texttt{Source branch} kan je kies je een branch op te kopiëren (momenteel is enkel \texttt{main} beschikbaar).
		\item Klik op de branch \texttt{vakken}, \texttt{Add file}, \texttt{Create new file}.
		\item Geef het nieuwe bestand de naam \texttt{vakken.txt} en zet er de vakken in die je dit semester volgt.
		\item Commit de wijziging op \github.
		\begin{itemize}
			\item Lokaal staan er momenteel twee branches: \texttt{main} en \texttt{hobbies}.
			\item Op \github\ staan ook twee branches:  \texttt{main} en \texttt{vakken}.
			\item Van de branch \texttt{hobbies} ontbreekt elk spoor op \github\ en het bestaan van de branch \texttt{vakken} is onbekend in \vscode.
		\end{itemize}
	\end{itemize}
	
	In wat volgt, proberen we alle branches samen te brengen. Dit kan uiteraard uitsluitend in \vscode, want \github\ heeft geen toegang tot je lokale systeem.
	
	\subsection{Fetch -- online informatie ophalen} 
	
	\begin{itemize}
		\item Klik op de branch-naam in \vscode, links onderaan in de statusbalk. In het lijstje bekende branches staan de lokale branches (\texttt{main} en \texttt{hobbies}) en de enige gekende remote branch (\texttt{origin/main}).
		\item Klik op \texttt{Fetch From All Remotes} naast \texttt{Graph}. Hierdoor gaat \git\ na wat er online is veranderd. Er worden op dit moment nog geen bestanden aangepast.
		\begin{itemize}
			\item De remote branch (\texttt{origin/vakken}) verschijnt nu ook in het lijstje bekende branches als je op de branch-naam in de statusbalk klikt.
		\end{itemize}
	\end{itemize}
	
	\subsection{Pull -- Branches van \github\ halen} 
	
	\begin{itemize}
		\item Klik in \vscode\ op de branch-naam in de statusbalk en in het lijstje dat verschijnt op \texttt{Create New Local Branch From...}.
		\item Kies voor de branch \texttt{origin/vakken}.
		\item Geef de nieuwe lokale branch de naam \texttt{vakken}.
		\begin{itemize}
			\item Verifieer dat nu zowel \texttt{origin/vakken} als \texttt{vakken} in de lijst bekende branches staan.
			\item Schakel je over naar de branch \texttt{vakken}, dan verschijnt het bestand \texttt{vakken.txt}.
			\item Het bestand \texttt{vakken.txt} verschijnt niet in de branch \texttt{hobbies} en het bestand \texttt{hobbies.txt} niet in de branch \texttt{vakken}.
		\end{itemize}
	\end{itemize}
	
	\subsection{Push -- Branches publiceren op \github} 
	
	\begin{itemize}
		\item Schakel over naar de branch \texttt{hobbies}. 
		\item Klik nu op \texttt{Publish Branch}, de branch wordt nu ook naar \github\ verzonden. 
		\begin{itemize}
			\item Op \github\ zijn nu ook de drie branches zichtbaar: \texttt{main}, \texttt{hobbies}, \texttt{vakken}.
			\item In het \texttt{Graph}-venster \vscode\ worden de drie branches ten opzicht van elkaar getoond, in verschillende kleuren. Naast elke branch-naam staat nu ook een cloud-symbooltje in dezelfde kleur omdat local en remote repository volledig gesynchroniseerd zijn. 
		\end{itemize}
	\end{itemize} 
	
	
	\subsection{Terminologie}
	
	\begin{itemize}
		\item \emph{Branch}: een alternatieve versie van de repository.
		\item \texttt{main}: de hoofdbranch\footnote{In principe kan de hoofdbranch een andere naam hebben, \texttt{master} of \texttt{trunk} komen ook frequent voor.} van de repository, waarin de stabiele versie van het project wordt bewaard.
		\item \emph{Fetch}: nagaan welke nieuwe branches en commits op \github\ beschikbaar zijn, zonder de lokale bestanden te wijzigen
		\item \emph{Pull} en \emph{Push}: wijzigingen worden afzonderlijk voor elke branch opgehaald of gepubliceerd.
	\end{itemize}
	
	Een typische workflow is om commits niet rechtstreeks op de hoofdbranch te maken, maar eerst een nieuwe branch te maken, daar de wijzigingen uit te proberen en pas als alles correct werkt, de wijzigingen samen te voegen met de hoofdbranch. Dat proces van samenvoegen heet \emph{merge} en komt in de volgende sectie aan bod.
	
	\section{Samenvoegen van branches: merge}
	\label{sec:merge}
	
	We zullen in dit deel drie soorten merges uitvoeren: die van nieuwe bestanden, van wijzigingen op verschillende plaatsen in hetzelfde bestand, en van wijzigingen op dezelfde plaats in een bestand, de gevreesde \emph{merge met conflicten}.
	
	\subsection{Lokale merge van nieuwe bestanden} Momenteel bevatten de branches \texttt{hobbies} en \texttt{vakken} elk een ander bestand dat nog niet in de branch \texttt{main} staat. We zullen beide branches samenvoegen met \texttt{main}, zodat alle bestanden er zijn samengebracht.
	
	\begin{itemize}
		\item Schakel over naar de \emph{doelbranch} \texttt{main}, de branch waarin aanpassingen zullen worden overgenomen.
		\item Klik in \texttt{Source Control}, naast \texttt{Changes} op \texttt{More Actions}, \texttt{Branch}, \texttt{Merge}.
		\item Kies in de lijst voor de \emph{bronbranch} \texttt{hobbies}, de branch waaruit aanpassingen zullen worden overgenomen.
		\begin{itemize}
			\item In \texttt{Explorer} verschijnt het bestand \texttt{hobbies.txt}.
			\item In de \texttt{Graph} komen de paden van \texttt{main} en \texttt{hobbies} opnieuw samen.
			\item Een merge resulteert eigenlijk in een commit van de aanpassingen uit de bronbranch in de doelbranch.
			\item De merge is asymmetrisch, de bronbranch verandert niet.
			\item De cloud-icoontjes tonen dat de merge in de remote repository nog niet is gebeurd.
		\end{itemize}
		\item Klik op de knop \texttt{Push} naast \texttt{Graph} om de lokale aanpassingen naar \github\ te pushen.
		\begin{itemize}
			\item In de \texttt{Graph} horen de vier branches \texttt{main}, \texttt{origin/main}, \texttt{hobbies} en \texttt{origin/hobbies} nu bij dezelfde node.
		\end{itemize}
	\end{itemize}
	
	Na een merge blijven beide branches bestaan. Onmiddellijk na de merge bevatten ze dezelfde wijzigingen, maar nadien kunnen ze opnieuw onafhankelijk van elkaar evolueren. In de praktijk gebeurt verdere ontwikkeling op een secundaire branch, waarna de aanpassingen later opnieuw kunnen worden samengevoegd met de hoofdbranch \texttt{main}.
	
	\subsection{Remote merge van nieuwe bestanden} De branch \texttt{\texttt{vakken}} zullen we op \github\ samenvoegen met de hoofdbranch. Dit simuleert de typische situatie waarbij er twee personen betrokken zijn bij een merge, de contributor en de reviewer. De eerste voert een aantal aanpassingen door in een eigen branch, de tweede gaat na of de aanpassingen geschikt zijn om in de hoofdbranch te worden toegevoegd.
	
	\begin{itemize}
		\item[] Een \emph{contributor} doet een bijdrage aan een repository in een eigen branch. Dit gebeurt typisch op het eigen lokale systeem, waarna de aanpassingen gepusht worden naar de \github-repository. Hier pasten we ter illustratie een online bestand rechtstreeks aan. Nu stelt de contributor voor om de aanpassingen in de hoofdbranch op te nemen.
		\item Klik in \github\ op \texttt{Pull requests} en vervolgens op \texttt{New pull request}.
		\item Kies in de drop-down list de ontvangende ``\texttt{base}''-branch \texttt{main} en rechts de branch \texttt{vakken} waaruit de aanpassingen zullen worden gehaald.
		\begin{itemize}
			\item Er verschijnt een lijst van alle commits die er ten opzichte van de \texttt{base}-branch zijn geweest.
			\item Daaronder staat een lijst van alle bestanden (of bestandsinhouden, zie verder) die zijn gewijzigd.
		\end{itemize}
		\item Klik op \texttt{Create pull request}, geef een beschrijving en klik opnieuw op \texttt{Create pull request}.
		\item[] In deze simulatie speel je beide rollen, maar in de praktijk zijn de volgende stappen voor de \emph{reviewer}, die een overzicht te zien krijgt van commits en eventuele conflicten (momenteel geen) en beslist om de merge al dan niet te laten doorgaan. 
		\item Klik op \texttt{Merge pull request}, voeg een commit message toe en bevestig met \texttt{Confirm merge}.
		\begin{itemize}
			\item Op \github\ bevatten \texttt{main} en \texttt{vakken} nu dezelfde bestanden.
			\item In de lokale repository is de merge nog niet gebeurd.
		\end{itemize}
		\item Klik in \vscode\ naast \texttt{Graph} op de knop \texttt{Fetch From All Remotes}.
		\begin{itemize}
			\item De \texttt{Graph} toont dat \texttt{origin/main} een nieuwere commit bevat dan \texttt{main} en dat de branch \texttt{vakken} daarbij samengevoegd wordt.
		\end{itemize}
		\item Klik naast \texttt{Graph} op de knop \texttt{Pull}.
		\begin{itemize}
			\item In \texttt{Explorer} zijn nu alle bestanden samengebracht in de \texttt{main}-branch.
			\item In de \texttt{Graph} bevinden \texttt{main} en \texttt{origin/main} zich opnieuw op dezelfde commit.
		\end{itemize}
	\end{itemize}
	
	\subsection{Merge zonder conflicten} We zullen nu twee verschillende wijzigingen maken op verschillende plaatsen in hetzelfde bestand.
	
	\begin{itemize}
		\item Maak in de lokale repository een nieuwe branch \texttt{programmeertaal},  voeg \emph{aan het einde} van het bestand \texttt{gegevens.txt} een nieuwe lijn toe met je favoriete programmeertaal, sla op en commit de wijziging.
		\item Terug in de branch \texttt{main} voeg je in het bestand \texttt{gegevens.txt} \emph{boven je naam} je studentennummer toe. Opslaan en committen. De programmeertaal is in dit bestand nog niet te zien!
		\item Merge de branch \texttt{programmeertaal} met de hoofdbranch \texttt{main} (\texttt{ctrl+shift+P}, \texttt{Git: Merge...}).
		\begin{itemize}
			\item Zowel het studentennummer als de favoriete programmeertaal verschijnen nu in het bestand \texttt{gegevens.txt} in de \texttt{main}-branch.
		\end{itemize}
	\end{itemize}
	
	\git\ kon beide wijzigingen automatisch combineren omdat ze verschillende delen van hetzelfde bestand aanpasten. Dit werkt zolang de inhoud van eenzelfde lijn niet is aangepast.
	
	\subsection{Een merge met conflict} We creëren eerst een situatie dat tot een conflict zal leiden, en zullen het daarna proberen oplossen.
	
	\begin{itemize}
		\item Maak een nieuwe branch \texttt{conflict}, verander in het bestand \texttt{gegevens.txt} de woonplaats in je werkelijke woonplaats, sla op en commit de wijziging.
		\item Terug in de branch \texttt{main} staat in \texttt{gegevens.txt} nog \texttt{Woonplaats: Gallifrey}. Verander dit in \texttt{Woonplaats: Camelot}. Dit simuleert een andere contributor die simultaan een (kennelijk opnieuw foute) aanpassing maakt. Opslaan en committen.
		\item Merge de branch \texttt{conflict } met \texttt{main} (\texttt{ctrl+shift+P}, \texttt{Git: Merge...}).
		\begin{itemize}
			\item \git\ detecteert conflicterende aanpassingen en kan nu niet automatisch beslissen welke versie correct is.
			\item Bekijk beide versies grondig in de editor en beslis welke actie er moet worden ondernomen..
		\end{itemize}
		\item Kies om enkel de lijn uit de doelbranch \texttt{main} te behouden, deze lijn te vervangen volgens wat in bronbranch staat, of om beide lijnen onder elkaar zetten.
		\item Pas na het afhandelen van het conflict, kan een commit uitgevoerd worden om de merge definitief te voltooien.
		\item Push nog de laatste aanpassingen naar \github\ (\texttt{Sync Changes}).
	\end{itemize}
	
	\subsection{Terminologie}
	
	\begin{itemize}
		
		\item \emph{Merge commit}: een commit die het resultaat van een merge bewaart.
		\item \emph{Bronbranch}: de branch waarvan de wijzigingen worden overgenomen.
		\item \emph{Doelbranch}: de branch waarin de wijzigingen worden opgenomen (hierboven steeds \texttt{main}).
		\item \emph{Pull request}: vraag van \emph{contributor} om aanpassingen in doelbranch op te nemen.
		\item De \emph{reviewer} inspecteert de wijzigingen vooraleer de merge uit te voeren.
		\item \emph{Merge conflict}: een situatie waarin \git\ niet automatisch kan bepalen welke versie moet worden behouden.
	\end{itemize}
	
\section{Starten van een bestaande repository}

Hierboven werd een repository lokaal aangemaakt, startend van een lege map, en vervolgens naar \github\ gepubliceerd. In de praktijk zal je echter vaak starten van een bestaande repository die zich al op GitHub bevindt: omdat je wil verder werken aan een reeds bestaande repository of omdat een teamlid de repository heeft geïnitialiseerd.

\subsection{Fork -- kopie van een repository van een andere eigenaar}\label{ssec:fork}

Zelfs wanneer een repository op \github\ publiek beschikbaar is, kan een gebruiker die geen expliciete schrijfrechten heeft gekregen, de repository enkel bekijken en kopiëren, maar niet rechtstreeks committen of wijzigingen mergen.

Om toch aan zo'n repository te kunnen bijdragen, gebruikt GitHub zogenaamde \emph{forks}. Een fork is een volledige kopie van een repository in een andere GitHub-account. De eigenaar van de oorspronkelijke repository behoudt de controle over de hoofdversie, terwijl de eigenaar van de fork vrij kan experimenteren in de eigen kopie.

\begin{itemize}
    \item Open de repository \href{https://github.com/srebry/test}{\texttt{github.com/srebry/test}}.
    \item Klik rechtsboven op \texttt{Fork}.
    \item Kies je eigen \github-account als bestemming.
    \item Bevestig.
    \begin{itemize}
        \item In je eigen account verschijnt nu de nieuwe repository. Interessant om weten: forks van publieke repositories blijven automatisch publiek.
        \item Alle bestanden, branches en commits werden gekopieerd.
		\item De nieuwe repository, de \emph{fork}, blijft verwijzen naar de bron, de zogenaamde \emph{upstream repository}, al kan je zelf geen merg uitvoeren naar deze bron.
		\item Net zoals bij een branch kunnen wijzigingen via een \emph{pull request} worden voorgesteld aan de upstream repository.
		\item Verschijnen er nieuwe commits in de upstream repository, dan kunnen die worden overgenomen door in \github\ te klikken op \texttt{Sync fork}, \texttt{Update branch}.
    \end{itemize}
\end{itemize}

Branches en forks lijken dus op elkaar: een branch is een alternatieve versie binnen dezelfde repository met dezelfde eigenaar. Een fork is een volledige kopie van een repository, typisch bedoeld voor een gebruiker die geen schrijfrechten heeft op de oorspronkelijke repository.

\subsection{Clone -- een online repository ophalen}

Het lokaal kopiëren van een online repository heet \emph{clonen}.

\begin{itemize}
    \item Open de fork die je zonet hebt gemaakt op \github.
    \item Klik op de groene knop \texttt{Code}.
    \item Kopieer de HTTPS-link van de repository.
    \item Om een nieuw project te beginnen, open je een nieuw venster van \vscode (\texttt{ctrl+shift+N}).
    \item Ga naar \texttt{Source Control} en klik op \texttt{Clone Repository}
    \item Plak de gekopieerde URL.
    \item Maak als doelmap een nieuwe map \texttt{test} in de map \texttt{C:\textbackslash Projecten}.
    \begin{itemize}
        \item Alle bestanden uit de repository verschijnen in \texttt{Explorer}.
        \item De volledige geschiedenis van commits en branches wordt mee gekopieerd en getoond in de \texttt{Graph}.
        \item Wijzigingen kunnen tussen de clone en de origin worden uitgewisseld met \texttt{Push} en \texttt{Pull}.
    \end{itemize}
\end{itemize}

Een clone is dus veel meer dan een gewone kopie van de bestanden die je via \texttt{Download ZIP} zou verkrijgen. De volledige repositorygeschiedenis wordt lokaal beschikbaar gemaakt én het blijft mogelijk om wijzigingen naar de remote repository te pushen.

\subsection{Terminologie}

\begin{itemize}
    \item \emph{Fork}: een kopie van een GitHub-repository in een andere GitHub-account.
    \item \emph{Upstream repository}: de oorspronkelijke repository waarvan een fork werd gemaakt.
    \item \emph{Clone}: een lokale kopie van een repository op de eigen computer.
    \item \emph{Origin}: de remote repository waarvan een clone afkomstig is.
\end{itemize}


	\section{Werkverdeling: issues en assignees}
	\label{sec:issues}
	
	Tot nu toe hebben we wijzigingen rechtstreeks uitgevoerd op bestanden. Bij grotere projecten is het echter nuttig om eerst vast te leggen \emph{wat} er moet gebeuren. Daarvoor gebruikt \github\ \emph{issues}. Deze beschrijven een taak die nog moet worden uitgevoerd, een nieuwe functie die moet worden toegevoegd of een fout die moet worden opgelost. Om er voor te zorgen dat niet iedereen hetzelfde probleem aan het oplossen is, kunnen issues toegekend worden aan een specifieke contributor, de \emph{assignee}.
	
	\subsection{Een issue en branch aanmaken}  We zullen een eerste issue aanmaken voor onze repository. Zoals steeds worden wijzigingen niet rechtstreeks op \texttt{main} uitgevoerd, maar in een nieuwe branch die specifiek aan de issue gekoppeld is. In de praktijk worden de taken in deze workflow over twee personen verdeeld, de reviewer en een contributor.
	
	\begin{itemize}
		\item[] De reviewer wil dat een contributor aan het bestand \texttt{gegevens.txt} de naam van de favoriete film toevoegt en voert volgende stappen uit.
		\item Open de repository op \github\ en klik bovenaan op \texttt{Issues}.
		\item Klik op \texttt{New Issue} en geef een titel (\texttt{1-film}) en beschrijving aan de issue.
		\item Klik op \texttt{Create}.
		\begin{itemize}
			\item De issue verschijnt nu in de lijst met openstaande issues.
		\end{itemize}
		\item Open de issue dat net is aangemaakt.
		\item Klik op de tekst \texttt{Create branch} in de rechter kolom onder \texttt{Development}.
		\item Geef een duidelijke branchnaam \texttt{1-film} die issue-nummer en omschrijving omvat.
		\begin{itemize}
			\item De branch bestaat momenteel enkel op \github.
			\item \github\ voorziet \git-commando's om de nieuwe branch naar de lokale repository te pullen, maar de werkwijze vanuit \vscode\ zoals hierboven werkt even goed.
		\end{itemize}
	\end{itemize}
	
	\subsection{Een assignee toekennen} Een issue kan aan één of meerdere personen worden toegewezen.
	
	\begin{itemize}
		\item Open opnieuw het zonet aangemaakte issue.
		\item Klik op het tandwiel naast \texttt{Assignees} in de rechter kolom.
		\item Selecteer een of meerdere contributors van de repository (voorlopig is dat enkel jijzelf).
		\begin{itemize}
			\item Het issues-overzicht toont op elk moment welke taken nog open staan.
			\item In de meest rechtse kolom zie je wie verantwoordelijk is voor een taak en welke taken nog niet zijn toegekend.
			\item De knop \texttt{assignees} laat toe de taken van een bepaalde contributor (bijvoorbeeld jezelf) te filteren.
		\end{itemize}
	\end{itemize}
	
	\subsection{De wijziging uitvoeren} Deze kleine aanpassing kunnen we rechtstreeks op \github\ doorvoeren. Grotere issues vergen de fetch-pull-commit-push-cyclus uit deel \ref{sec:branches}.
	
	\begin{itemize}
		\item[] De contributor die de issue zal opvolgen, voert volgende stappen uit.
		\item Open op \github\ in de branch \texttt{1-film} het bestand \texttt{gegevens.txt}. Voeg een nieuwe lijn toe met daarop de naam van je favoriete film en sla het bestand op.
		\item Maak een commit. In de \texttt{1-film}-branch is de gevraagde issue opgelost, maar nog niet in de \texttt{main}-branch.
		\item In dit geval lost één commit de issue op, de contributor doet een pull request.
	\end{itemize}
	
	
	\subsection{De issue sluiten} In de volgende stap,  zal de reviewer moeten nagaan of de aanpassingen de issue effectief oplossen.
	
	\begin{itemize}
		\item De reviewer behandelt de pull request en klikt op \texttt{Confirm merge} als de aanpassingen mogen worden overgenomen in de doelbranch en de issue daarmee opgelost is.
		\begin{itemize}
			\item Doordat de branch \texttt{1-film} aan de issue is gekoppeld, sluit deze merge de issue automatisch af. De issue is nog terug te vinden bij de \texttt{Closed issues}
			\item Blijkt de merge toch niet afdoende, dan kan de issue opnieuw worden geopend: \github\, \texttt{Issues}, \texttt{Closed}, klik op het onterecht gesloten issue, \texttt{Reopen issue}.
			\item Is de issue op een andere manier opgelost, zonder merge van de aangehechte branch, dan kan deze manueel worden afgesloten: \github\, \texttt{Issues}, klik op de issue die kan worden afgesloten, \texttt{Close issue}.
			\item In het closed issue overzicht is het mogelijk om na te gaan wie welke deeltaak heeft opgelost en dus aan wie vragen daarover te richten.
		\end{itemize}
	\end{itemize}
	
	\subsection{Terminologie}
	
	\begin{itemize}
		\item \emph{Issue}: een probleem, taak of vraag die moet worden opgevolgd. Issues kunnen erg uiteenlopend zijn, ze kunnen een fout (\emph{bug}) in een bestaand document signaleren, of vragen om een extra functie (\emph{feature request}).
		\item \emph{Assignee}: de contributor die is toegekend aan het oplossen van een specifiek issue.
		\item \emph{Open issue}: een issue waarvoor nog geen oplossing beschikbaar is.
		\item \emph{Closed issue}: een issue dat werd opgelost of afgewerkt.
	\end{itemize}
	
	\section{De teamopdacht, repository en workflows}
	\label{sec:teamopdracht}
	
	\subsection{Team repository} Elk team werkt in een fork van de repository \texttt{teamopdrachtX} met \texttt{X} de letter het team. De teamleider zal de verantwoordelijkheid hebben om alle eindproducten (opgeloste oefeningen, Python-scripts, eindverslag en poster, \ldots) aan de \texttt{main}-branch van deze repository toe te voegen door verschillende branches te mergen.
	
	\subsection{Wekelijkse deelopdrachten} De wekelijkse opgave wordt in de map \texttt{0-bronnen} van de upstream repository geplaatst door de begeleiders. De knop \texttt{Sync fork} in \github\ laat toe om deze naar de team-repository te halen.

	\subsection{Vergaderverslagen} De vergaderverslagen worden door de verslaggever voorbereid in de branch \texttt{1-vergaderingen}. Ze worden pas in de \texttt{main}-branch gemerged door de teamleider nadat ze zijn goedgekeurd door het team (meestal pas de volgende vergadering).
	
	\subsection{Wekelijkse branches} Voor de deelopdrachten zal telkens een nieuwe branch gemaakt worden, startend van de hoofdbranch \texttt{main}. De branch \texttt{2-oefeningen} is al voorzien, de komende weken komen daar onder meer \texttt{implementatie}, \texttt{verslag} en \texttt{poster} bij, horend bij de deeltaak die jullie in de betreffende week krijgen. De taakverantwoordelijke zal de verantwoordelijkheid hebben om alle deelopdrachten (afzonderlijke oefeningen, scripts of bestanden) in deze branch samen te brengen.
	
	\subsection{Issues voor alle deeltaken} Elke week overlegt het team welke deelopdrachten er zijn en de taakverantwoordelijke maakt de corresponderende issues aan. Bij het deel oefeningen zijn dat eenvoudigweg de te maken oefeningen, issues \texttt{oefening01} tot \texttt{oefening13}.
	
	Uiteraard is het mogelijk dat er na verloop van tijd nog meer issues bijkomen, bijvoorbeeld een bug die ontdekt wordt in een Python-script of een stuk tekst dat moet herschreven worden.
	
	\subsection{Taakverdeling met assignees} Tijdens de teamvergadering wordt beslist hóe jullie de taken zullen verdelen. Om het overzicht te bewaren is het aangewezen dat de issues niet allemaal tegelijk worden verdeeld over de teamleden, maar dat elk teamlid maar één issues markeert waar die op dat moment effectief mee bezig is.
	
	\subsection{Feature branches en langlevende branches}

	Er blijkt dus dat er in de teamopdracht, net zoals in het merendeel van de \git-repositories, twee soorten branches zullen worden gebruikt. In de voorbeelden van de eerste secties van deze tekst, werden branches gebruikt om een beperkte wijziging uit te werken. Nadat de wijziging was gemerged met de hoofdbranch, had de branch haar doel bereikt. Dergelijke branches noemt men \emph{feature branches}. Na de merge wordt een feature branch vaak verwijderd omdat alle wijzigingen ondertussen in de hoofdbranch aanwezig zijn. Voorbeelden in de teamopdracht zijn\ldots
\begin{itemize}
\item \texttt{oefening01}, \ldots, \texttt{oefening13}
\item \texttt{bug-python-script}
\item \texttt{spellingscontrole-hoofdstuk5}
\item \texttt{afbeelding-maken-voor-poster}
\end{itemize}

Niet alle branches zijn echter tijdelijk. Sommige projecten gebruiken meerdere branches die permanent blijven bestaan en elk een eigen rol hebben. Men spreekt dan van \emph{langlevende branches}. Voorbeelden hiervan in de context van de teamopdracht zijn
\begin{itemize}
\item \texttt{teamX-main}
\item \texttt{teamX-0-vergaderingen}
\item \texttt{teamX-0-oefeningen}
\end{itemize}

	\subsection{Commit, push en merge}
	
	Tracht gestructureerd te werken, één deelprobleempje tegelijk, en voer snel genoeg een commit uit. Die documenteert dan achteraf wat er al is gebeurd.
	
	Zolang je eigen bijdragen niet naar \github\ zijn gepushed, bestaan deze niet voor andere teamleden en de begeleiding. Aarzel niet om dat af en toe te doen:
	\begin{itemize}
		\item Onmiddellijk bij het afwerken van een issue.
		\item Blijf niet urenlang staren naar een opdracht als je vast komt te zitten. Beschrijf het probleem waar je mee worstelt en deel het als een nieuwe issue of sub-issue.
		\item Push ál je materiaal aan het einde van de teamsessie.
	\end{itemize}

	\section*{Besluit}

\git\ en \github\ vormen vandaag de standaard voor versiebeheer en samenwerking in professionele projecten. Ze worden gebruikt van zeer kleine projecten door één gebruiker, tot grote software-, onderzoeks- en ingenieursprojecten waaraan honderden of zelfs duizenden personen tegelijk bijdragen.

Bij een eerste kennismaking lijkt \git\ vaak complex. De begrippenlijst is ellenlang -- \emph{repository}, \emph{commit}, \emph{branch}, \emph{merge}, \emph{issue}, \emph{pull request}, \emph{contributor} en \emph{reviewer} -- en deze uitgebreide terminologie is niet altijd intuïtief. Die begrippen zijn echter noodzakelijk om de onderliggende ideeën te begrijpen en om efficiënt te kunnen samenwerken in grotere projecten.

Door \git\ en \github\ consequent te gebruiken ontstaan verschillende voordelen. Dankzij commits en repositories gaat informatie niet verloren wanneer bestanden per ongeluk worden verwijderd of onbedoeld worden aangepast.  Nieuwe ideeën kunnen eerst worden ontwikkeld in afzonderlijke branches zonder de stabiele versie van het project in gevaar te brengen en meerdere personen kunnen gelijktijdig aan verschillende onderdelen werken. Tot slot biedt werken met \git\ en \github\ niet alleen technische hulpmiddelen voor versiebeheer, maar ondersteunt ook de organisatie van het project zelf. Issues, assignees, branches en pull requests maken het teamwerk zichtbaar.

\git\ en \github\ zijn dus veel meer dan een systeem om bestanden op te slaan. Ze vormen een hulpmiddel om kennis, documenten, code en verantwoordelijkheden op een gestructureerde manier te beheren.


\end{document}